\documentclass[11pt]{article}
\usepackage[T1]{fontenc}
\usepackage[margin=0.88in]{geometry}
\usepackage{amsmath,amssymb}
\usepackage{booktabs,tabularx}
\usepackage[hidelinks]{hyperref}
\usepackage{microtype}

\newcommand{\R}{\mathbb R}
\newcommand{\Sph}{\mathbb S}
\newcommand{\LC}{\mathrm{LC}}

\title{The two-measure functional Minkowski problem}
\author{Literature-already-answered note for arXiv:2105.09168}
\date{11 August 2026}

\begin{document}
\maketitle

\paragraph{Status.}
\textbf{Completely answered by Falah--Rotem, arXiv:2502.16929,
Theorem~1.5.}
The later paper gives necessary and sufficient conditions for a pair of
measures to be the two surface-area measures of an integrable log-concave
function, and proves uniqueness up to translation.  This is precisely
Problem~6.7 of the source paper, not merely a related one-measure or smooth
special case.

\section*{The source question}
For an upper-semicontinuous integrable log-concave function
$g=e^{-\phi}$ on $\R^n$, the source associates two finite measures:
\[
 \mu_g=(\nabla\phi)_{\#}(g\,dx)\quad\hbox{on }\R^n,
 \qquad
 \nu_g=(n_{\operatorname{supp}g})_{\#}
 \bigl(g\,\mathcal H^{n-1}|_{\partial\operatorname{supp}g}\bigr)
 \quad\hbox{on }\Sph^{n-1}.
\]
After treating the essentially continuous case $\nu_g=0$, Problem~6.7 on
source PDF page~21 asks:
\begin{quote}
Fix a finite Borel measure $\mu$ on $\R^n$ and a finite Borel measure
$\nu$ on $\Sph^{n-1}$.  Under what conditions can one find
$g\in\LC_n$ with $\mu_g=\mu$ and $\nu_g=\nu$?
\end{quote}
The surrounding text says that solving this problem is needed to extend the
Riesz-type first-variation representation beyond the essentially continuous
class.

\section*{The exact later theorem}
Theorem~1.5 of arXiv:2502.16929 (supporting PDF page~5; source label
\texttt{thm:main-Minkowksi}) states the following.  Let $\mu$ be a finite
Borel measure on $\R^n$ and $\nu$ a finite Borel measure on
$\Sph^{n-1}$.  There is an integrable log-concave function $f$ such that
$(\mu_f,\nu_f)=(\mu,\nu)$ if and only if:
\begin{enumerate}
\item $\mu\ne0$;
\item $\mu$ has finite first moment and the pair is centered, meaning
\[
 \int_{\R^n}x\,d\mu(x)+
 \int_{\Sph^{n-1}}\theta\,d\nu(\theta)=0;
\]
\item $\mu$ and $\nu$ are not supported on a common hyperplane
$H\subseteq\R^n$.
\end{enumerate}
Moreover, the realizing function is unique up to translation: if $f$ and
$g$ have the same pair, then $g(x)=f(x+v)$ for some $v\in\R^n$.

The theorem's ``common hyperplane'' condition uses the sphere as a subset of
$\R^n$: there is no hyperplane $H\subseteq\R^n$ supporting both measures.

\section*{Why the identification is exact}
\begin{center}
\small
\begin{tabularx}{\textwidth}{@{}p{0.29\textwidth}XX@{}}
\toprule
Feature & Problem~6.7 & Falah--Rotem Theorem~1.5 \\
\midrule
Interior measure & $\mu_g$ on $\R^n$ & the same $\mu_f$ \\
Boundary measure & $\nu_g$ on $\Sph^{n-1}$ & the same $\nu_f$ \\
Function class & integrable log-concave functions & the same class \\
Requested conclusion & characterize existence & iff conditions, plus uniqueness \\
\bottomrule
\end{tabularx}
\end{center}

The 2025 paper's abstract also says explicitly that it ``fully solve[s]'' the
functional Minkowski problem.  Its theorem includes the earlier moment-measure
case $\nu=0$ but is not restricted to it.  Thus the source problem is no
longer open as stated.

\paragraph{Novelty status.}
This packet records an exact later-literature resolution.  It does not claim
a new theorem or an independent proof from this run.

\paragraph{Human review.}
Check Problem~6.7 on source PDF page~21 and Theorem~1.5 on supporting PDF
page~5.  The proof of the existence direction begins in Section~3 of the
supporting paper; uniqueness is cited there as previously known in full
generality.

\begin{thebibliography}{9}
\raggedright
\bibitem{source}
L. Rotem,
\emph{A Riesz representation theorem for log-concave functions},
arXiv:2105.09168.  See Problem~6.7 on PDF page~21.

\bibitem{answer}
T. Falah and L. Rotem,
\emph{On the functional Minkowski problem},
arXiv:2502.16929.  See Theorem~1.5 on PDF page~5.
\end{thebibliography}

\end{document}
